% ============================================================
% BÖLÜM 3: ARAŞTIRMANIN AMACI VE ALT AMAÇLARI
% ============================================================

\chapter{ARAŞTIRMANIN AMACI VE ALT AMAÇLARI}

\section{Genel Amaç}

Bu çalışmanın genel amacı, okul psikolojik danışmanlarının House-Tree-Person (HTP) çizimlerini değerlendirme süreçlerini desteklemek üzere, yapay zekâ tabanlı bir karar destek sistemi geliştirmek ve bu sistemin PDR alanındaki potansiyel katkılarını ortaya koymaktır.

Karar destek sistemleri (Decision Support Systems - DSS), karmaşık problem çözme süreçlerinde profesyonellere yardımcı olan bilgi sistemleridir (Power, 2002). Sağlık alanında yaygın biçimde kullanılan klinik karar destek sistemleri, tanı doğruluğunu artırmakta ve tedavi süreçlerini optimize etmektedir (Sutton ve ark., 2020). NIDA, bu yaklaşımı PDR alanına uyarlamaktadır.

\section{Alt Amaçlar}

Çalışmanın genel amacı doğrultusunda aşağıdaki alt amaçlar belirlenmiştir:

\begin{enumerate}
  \item \textbf{Akademik Bilgi Tabanı Oluşturma:}
  \begin{itemize}
    \item HTP testinin temel kaynaklarını (Buck, 1948; Hammer, 1958; Machover, 1949; Koch, 1952) sistematik biçimde derlemek
    \item Güncel akademik literatürü tarayarak yeni bulguları entegre etmek
    \item Çelişkili akademik görüşleri kayıt altına almak ve şeffaf biçimde sunmak
    \item Yorumlama kurallarını Lowenfeld (1947) ve Koppitz (1968) gelişim evrelerine göre farklılaştırmak
    \item Kanıt gücüne dayalı (evidence-based) bir puanlama sistemi geliştirmek
  \end{itemize}

  \item \textbf{Teknolojik Altyapı Geliştirme:}
  \begin{itemize}
    \item Web tabanlı, erişilebilir bir kullanıcı arayüzü tasarlamak (WCAG 2.1 standartlarına uygun)
    \item Çok modlu dil modeli (Gemini 2.0 Flash) entegrasyonu sağlamak
    \item Güvenli veri depolama ve kullanıcı yönetimi sistemi kurmak (KVKK uyumlu)
    \item Raporlama ve analiz geçmişi fonksiyonları geliştirmek
    \item Yüksek erişilebilirlik için responsive tasarım uygulamak
  \end{itemize}

  \item \textbf{Etik Çerçeve Belirleme:}
  \begin{itemize}
    \item Sistemin sınırlarını APA (2017) ve TPD (2004) etik ilkelerine uygun biçimde tanımlamak
    \item Kullanıcı uyarıları ve sorumluluk reddi metinleri oluşturmak
    \item Çocuk verilerinin korunmasına ilişkin politikalar belirlemek (BM Çocuk Hakları Sözleşmesi uyumu)
    \item PDR etik ilkeleri ile uyumu sağlamak
    \item Yapay zekâ etiği ilkelerini (UNESCO, 2021) göz önünde bulundurmak
  \end{itemize}

  \item \textbf{Kullanılabilirlik Sağlama:}
  \begin{itemize}
    \item Okul psikolojik danışmanlarının ihtiyaçlarına uygun, sezgisel arayüz tasarlamak
    \item Türkçe ve İngilizce dil desteği sunmak
    \item Yorumları anlaşılır ve uygulanabilir biçimde sunmak
    \item Eğitici içeriklerle kullanıcı yetkinliğini desteklemek
    \item Mobil cihaz uyumluluğu sağlamak
  \end{itemize}
\end{enumerate}

\section{Araştırma Soruları}

Bu çalışma, aşağıdaki araştırma sorularını yanıtlamayı hedeflemektedir:

\begin{enumerate}
  \item HTP yorumlama literatürü, yapay zekâ tabanlı bir karar destek sistemi için yapılandırılmış bir bilgi tabanına nasıl dönüştürülebilir?
  \item Çok modlu büyük dil modelleri (MLLM), çocuk çizimlerinin analizi için nasıl optimize edilebilir?
  \item Akademik literatürdeki çelişkili yorumlar, kullanıcılara nasıl şeffaf biçimde sunulabilir?
  \item Yapay zekâ tabanlı değerlendirme araçları, PDR etik ilkeleri ile nasıl uyumlu hale getirilebilir?
\end{enumerate}

\newpage
